%%%%%%%%%%%%%%%%%%%%%%%%%%%%%%%%%%%%%%%%%%%%%%%%%%%%%%%%%%%%%%%%%%%%%%
% Lớp 11 | Giới hạn. Hàm số liên tục | Giới hạn của hàm số
% Dạng 1D3?2-6: PP lượng liên hợp, kết quả vô cực
% Số câu: 1  (sinh tự động bởi tach_cau.py; từ nay thầy sửa trực tiếp file này)
%%%%%%%%%%%%%%%%%%%%%%%%%%%%%%%%%%%%%%%%%%%%%%%%%%%%%%%%%%%%%%%%%%%%%%

%%%%%%%%%%%%%%%% TRẮC NGHIỆM 4 PHƯƠNG ÁN %%%%%%%%%%%%%%%%

% nguồn: Ok_Toan11_TracNghiem_Lan2
\begin{ex}%[1D3H2-6]%[ma:f9163587]
	Giới hạn  $\lim\limits_{x\to -\infty } \left( 3x-\sqrt{9x^2-1} \right)$ bằng
	\choice 
		{$+\infty $} 
		 {$0$}
		{\True $-\infty  $}
		 {$-1$}
	 \loigiai{
		Ta có \begin{align*}
		&\lim\limits_{x\to -\infty } \left( 3x-\sqrt{9x^2-1} \right) \\ 
		&= \lim\limits_{x\to -\infty  } \left( 3x+x\sqrt{9-\dfrac{1}{x^2}} \right) \\
		&= \lim\limits_{x\to  -\infty } x\left( 3+\sqrt{9-\dfrac{1}{x^2}} \right) \\
		&=  -\infty.
	\end{align*}
} 
\end{ex}

% nguồn: 2026-11-15_11_1D3_Gioi_han (nhập 14/11/2026)
\begin{bt}%[1D3H2-6]%[ma:29ff7f31]
% nguon: Archive/SP_SauLuoi_2024-2025/SP_HKI/1-TN-DS-TLN-THPT-NguyenThiMinhKhai-HCM-HKI-NH24-25.tex #9
%[KNTT - Lớp 11 - Ôn tập giữa học kì 1 - Đề Nguyễn Thị Minh Khai]%[Thanh Phát]
\begin{enumerate}
	\item Tính giới hạn $\lim\limits_{x\to +\infty} \left(\sqrt{4x^2+x}-2024x \right)$.
	\item Tính giới hạn $\lim\limits_{x\to \left(-2\right)^{-} } \dfrac{1-\sqrt{x+3}}{x^2+4x+4}$.
\end{enumerate}

\loigiai{
\begin{enumerate}
	\item Ta có $\lim\limits_{x\to +\infty} \left(\sqrt{4x^2+x}-2024x \right)=\lim\limits_{x\to +\infty}\left[ x\cdot \left( \sqrt{4+\dfrac{1}{x}}-2024 \right) \right]$.\\
	Vì $\heva{&\lim\limits_{x\to +\infty} x =+\infty\\
		&\lim\limits_{x\to +\infty} \left( \sqrt{4+\dfrac{1}{x}} -2024 \right)=-2022 <0 }$ nên $\lim\limits_{x\to +\infty}\left[ x\cdot \left( \sqrt{4+\dfrac{1}{x}}-2024 \right) \right] =-\infty$.\\
	Vậy $\lim\limits_{x\to +\infty} \left(\sqrt{4x^2+x}-2024x \right)=-\infty$.
	\item Ta có $\lim\limits_{x\to \left(-2\right)^{-}} \dfrac{1-\sqrt{x+3}}{x^2+4x+4} = \lim\limits_{x\to (-2)^-} \dfrac{-(x+2)}{(x+2)^2\left(1+\sqrt{x+3} \right)} = \lim\limits_{x\to (-2)^-} \left( \dfrac{1}{x+2}\cdot \dfrac{-1}{1+\sqrt{x+3}} \right)$.\\
	Vì $\heva{&\lim\limits_{x\to (-2)^-} \dfrac{1}{x+2} = -\infty\\ &\lim\limits_{x\to (-2)^-} \dfrac{-1}{1+\sqrt{x+3}}=-\dfrac{1}{2} }$ nên $\lim\limits_{x\to (-2)^-} \left( \dfrac{1}{x+2}\cdot \dfrac{-1}{1+\sqrt{x+3}} \right)=+\infty$.\\
	Vậy $\lim\limits_{x\to \left(-2\right)^{-}} \dfrac{1-\sqrt{x+3}}{x^2+4x+4} =+\infty$.
\end{enumerate}
}
\end{bt}

% nguồn: chuyen_dang (nhập 18/11/2026)
\begin{ex}%[1D3V2-6]%[ma:3b8f8864]
	Tìm giới hạn $A=\lim\limits_{x  \to+\infty}\left(\dfrac{1+x}{x-1}\right)^x$.
	\loigiai{Với $A=\lim\limits_{x \to+\infty}\left(\dfrac{1+x}{x-1}\right)^x$, đặt $t=x-1$, khi $x \to+\infty$ thì  $t \to+\infty$. Do đó
	$$ 
	A=\lim\limits_{t \to+\infty}\left(\dfrac{t+2}{t}\right)^{t+1}=\lim\limits_{t \to+\infty} \dfrac{t+2}{t}\left(\dfrac{t+2}{t}\right)^t=\lim\limits_{t \to+\infty}  \dfrac{t+2}{t}\cdot \left[\left(1+\dfrac{2}{t}\right)^{\tfrac{t}{2}}\right]^2=\mathrm{e}^2.
	 $$
	 Vậy $A=\mathrm{e}^2$.
	} 
\end{ex}

% nguồn: chuyen_dang (nhập 18/11/2026)
\begin{ex}%[1D3V2-6]%[ma:3b8f8864]
	Tìm giới hạn $A=\lim\limits_{x  \to+\infty}\left(\dfrac{1+x}{x-1}\right)^x$.
	\loigiai{Với $A=\lim\limits_{x \to+\infty}\left(\dfrac{1+x}{x-1}\right)^x$, đặt $t=x-1$, khi $x \to+\infty$ thì  $t \to+\infty$. Do đó
	$$ 
	A=\lim\limits_{t \to+\infty}\left(\dfrac{t+2}{t}\right)^{t+1}=\lim\limits_{t \to+\infty} \dfrac{t+2}{t}\left(\dfrac{t+2}{t}\right)^t=\lim\limits_{t \to+\infty}  \dfrac{t+2}{t}\cdot \left[\left(1+\dfrac{2}{t}\right)^{\tfrac{t}{2}}\right]^2=\mathrm{e}^2.
	 $$
	 Vậy $A=\mathrm{e}^2$.
	} 
\end{ex}
